Vzdělávací profil jazykových znalostí ukazuje silnou pozitivní korelaci mezi dosaženým vzděláním a~jazykovou vybaveností ve všech zemích EU27; pro ČR je charakteristické, že mezera mezi terciárně vzdělanými a~pracovníky s~výučním listem je výraznější než v~AT nebo DE.
Tato vzdělávací disproporce jazykových znalostí vytváří segmentaci i~na trhu práce: nízce vzdělaní pracovníci (typicky průmysloví dělníci ve zpracovatelství) jsou bez jazykové kompetence mnohem více závislí na domácích institucích ochrany mezd --- tedy na kolektivních smlouvách --- než jejich vysoce kvalifikovaní kolegové.
Právě tato skupina je přitom v~ČR nejméně chráněna: ve zpracovatelském průmyslu (\ac{NACE} C), kde jsou koncentrováni, není sektorová \ac{KS} s~erga omnes pokrytím standardem.
Sociální dialog v~ČR by se tedy měl zaměřit na skupiny s~nejnižší individuální vyjednávací silou --- nízce vzdělané dělníky s~omezenou jazykovou mobilitou --- jako na primární beneficienty kolektivní ochrany mezd.

Vzdělávací profil jazykových znalostí ukazuje silnou pozitivní korelaci mezi dosaženým vzděláním a~jazykovou vybaveností ve všech zemích EU27; pro ČR je charakteristické, že mezera mezi terciárně vzdělanými a~pracovníky s~výučním listem je výraznější než v~AT nebo DE.
Tato vzdělávací disproporce jazykových znalostí vytváří segmentaci i~na trhu práce: nízce vzdělaní pracovníci (typicky průmysloví dělníci ve zpracovatelství) jsou bez jazykové kompetence mnohem více závislí na domácích institucích ochrany mezd --- tedy na kolektivních smlouvách --- než jejich vysoce kvalifikovaní kolegové.
Právě tato skupina je přitom v~ČR nejméně chráněna: ve zpracovatelském průmyslu (\ac{NACE} C), kde jsou koncentrováni, není sektorová \ac{KS} s~erga omnes pokrytím standardem.
Sociální dialog v~ČR by se tedy měl zaměřit na skupiny s~nejnižší individuální vyjednávací silou --- nízce vzdělané dělníky s~omezenou jazykovou mobilitou --- jako na primární beneficienty kolektivní ochrany mezd.

Vzdělávací profil jazykových znalostí ukazuje silnou pozitivní korelaci mezi dosaženým vzděláním a~jazykovou vybaveností ve všech zemích EU27; pro ČR je charakteristické, že mezera mezi terciárně vzdělanými a~pracovníky s~výučním listem je výraznější než v~AT nebo DE.
Tato vzdělávací disproporce jazykových znalostí vytváří segmentaci i~na trhu práce: nízce vzdělaní pracovníci (typicky průmysloví dělníci ve zpracovatelství) jsou bez jazykové kompetence mnohem více závislí na domácích institucích ochrany mezd --- tedy na kolektivních smlouvách --- než jejich vysoce kvalifikovaní kolegové.
Právě tato skupina je přitom v~ČR nejméně chráněna: ve zpracovatelském průmyslu (\ac{NACE} C), kde jsou koncentrováni, není sektorová \ac{KS} s~erga omnes pokrytím standardem.
Sociální dialog v~ČR by se tedy měl zaměřit na skupiny s~nejnižší individuální vyjednávací silou --- nízce vzdělané dělníky s~omezenou jazykovou mobilitou --- jako na primární beneficienty kolektivní ochrany mezd.

Vzdělávací profil jazykových znalostí ukazuje silnou pozitivní korelaci mezi dosaženým vzděláním a~jazykovou vybaveností ve všech zemích EU27; pro ČR je charakteristické, že mezera mezi terciárně vzdělanými a~pracovníky s~výučním listem je výraznější než v~AT nebo DE.
Tato vzdělávací disproporce jazykových znalostí vytváří segmentaci i~na trhu práce: nízce vzdělaní pracovníci (typicky průmysloví dělníci ve zpracovatelství) jsou bez jazykové kompetence mnohem více závislí na domácích institucích ochrany mezd --- tedy na kolektivních smlouvách --- než jejich vysoce kvalifikovaní kolegové.
Právě tato skupina je přitom v~ČR nejméně chráněna: ve zpracovatelském průmyslu (\ac{NACE} C), kde jsou koncentrováni, není sektorová \ac{KS} s~erga omnes pokrytím standardem.
Sociální dialog v~ČR by se tedy měl zaměřit na skupiny s~nejnižší individuální vyjednávací silou --- nízce vzdělané dělníky s~omezenou jazykovou mobilitou --- jako na primární beneficienty kolektivní ochrany mezd.

\input{texparts/python/language_skills_isced_map}



